\documentclass[11pt,a4paper]{article}
\usepackage[T1]{fontenc}
\usepackage{lmodern}
\usepackage[margin=27mm]{geometry}
\usepackage{amsmath,amssymb,amsthm,mathtools}
\usepackage{microtype}
\usepackage{booktabs,array}
\usepackage{enumitem}
\usepackage[hidelinks]{hyperref}
\setlength{\parskip}{5pt plus 1pt}
\setlength{\parindent}{0pt}
\setlength{\emergencystretch}{3em}
\setlist{itemsep=2pt,topsep=3pt}
\raggedbottom
\newtheorem{theorem}{Theorem}[section]
\newtheorem{lemma}[theorem]{Lemma}
\newtheorem{corollary}[theorem]{Corollary}
\newtheorem{proposition}[theorem]{Proposition}
\theoremstyle{definition}
\newtheorem{definition}[theorem]{Definition}
\newtheorem{example}[theorem]{Example}
\theoremstyle{remark}
\newtheorem{remark}[theorem]{Remark}
\newcommand{\Z}{\mathbb{Z}}
\newcommand{\N}{\mathbb{N}}
\newcommand{\Q}{\mathbb{Q}}
\newcommand{\F}{\mathbb{F}}
\newcommand{\ps}[1]{#1[[x]]}
\newcommand{\coeff}[2]{[x^{#1}]\,#2}
\newcommand{\hl}[1]{\textbf{#1}}
\newcommand{\para}[1]{\par\medskip\noindent\textbf{#1}\enspace}
\hypersetup{pdftitle={Parity of an infinitely nested series reversion: a proof of a conjecture in OEIS A195194},pdfauthor={Dev D. Goyal},pdfsubject={OEIS A195194}}
\title{Parity of an infinitely nested series reversion\\[6pt]\large A proof of a conjecture in OEIS A195194}
\author{Dev D. Goyal}
\date{October 2026}
\begin{document}
\maketitle
\begin{abstract}
The sequence A195194 of the On-Line Encyclopedia of Integer Sequences (OEIS), contributed by Paul D.~Hanna, is defined
by an infinitely nested expression,
\[
A(x)=\mathrm{INV}\Bigl(x-x\,\mathrm{INV}\bigl(x-2x\,\mathrm{INV}(x-3x\,\mathrm{INV}(x-\cdots))\bigr)\Bigr),
\]
where $\mathrm{INV}$ denotes series reversion, the operation that undoes a substitution of power series. Its terms are
$1,1,4,29,312,4528,\dots$. In 2026 Hanna conjectured that $a(n)$ is odd exactly when $n$ is a power of $2$. We first give
the infinite nesting a precise meaning: it defines exactly one power series, and cutting it off after $d$ levels gives
the correct coefficients up to degree $d+1$. This holds for arbitrary weights in place of $1,2,3,\dots$. We then prove the
conjecture. Modulo $2$ the weight $2$ at the second level vanishes, the whole nest collapses, and $A$ satisfies the
Catalan equation $A=x+A^2$, whose solution modulo $2$ is $x+x^2+x^4+x^8+\cdots$. The results have been checked by
computer and proved formally in the Lean proof assistant.
\end{abstract}
\begin{center}\begin{minipage}{0.88\linewidth}\small
\textbf{How this paper was produced.} This paper reports the result of an autonomous AI research effort.
ChatGPT (OpenAI), using the GPT-6 model with Pro reasoning and working from prompts by Dev~D.~Goyal, found the proof, wrote its Lean~4 formalization, and wrote a first draft of the paper.
Dev~D.~Goyal chose the problems and asked for checks and revisions, but did not supply proof ideas.
Claude Opus~5.5, used with Claude Code, then checked the work independently: it rebuilt the formal proofs from
source and confirmed that they rely on no unproved assumptions, recomputed the sequence from its definition and tested
every claim against the published data, and read every step of the proofs. It also rewrote this paper for clarity.
Dev~D.~Goyal has verified the results and is responsible for them. The paper has not been peer reviewed.
The section ``About this paper'' at the end gives the details.
\end{minipage}\end{center}

\section{Introduction}

\subsection{The sequence}\label{sec:sequence}
The \emph{On-Line Encyclopedia of Integer Sequences} (OEIS)~\cite{oeisdb} is a large public database of integer
sequences, used by mathematicians and scientists to identify sequences and to record what is known about them. Many
entries also record \emph{conjectures}: patterns that the contributor observed in the data but could not prove.

This paper is about entry A195194~\cite{oeis}, contributed by Paul~D.~Hanna in 2011. The entry reads, in part:
\begin{quote}\small\ttfamily\raggedright G.f.: A(x) = INV(x - x*INV(x - 2*x*INV(x - 3*x*INV(x - 4*x*INV(x - ...))))), where INV(F(x)) = series reversion of F(x).\par\smallskip Conjecture: a(n) is odd iff n is a power of 2. - Paul D. Hanna, Apr 07 2026\par\smallskip\hfill\textrm{--- OEIS A195194, Paul D. Hanna, September 2011}\end{quote}

Here ``G.f.''\ stands for \emph{generating function}: the terms $a(1),a(2),a(3),\dots$ of the sequence are stored as the
coefficients of the power series $A(x)=a(1)\,x+a(2)\,x^2+a(3)\,x^3+\cdots$.

\emph{Series reversion} is the power-series version of an inverse function. If $F(x)=x+f_2x^2+f_3x^3+\cdots$, then
$\mathrm{INV}(F)$ is the series $V(x)=x+v_2x^2+\cdots$ for which substituting $V$ into $F$ gives back $x$:
$F(V(x))=x$. For example, $\mathrm{INV}(x-x^2)$ is the series $x+x^2+2x^3+5x^4+14x^5+\cdots$ of the Catalan numbers.
Section~\ref{sec:background} explains this in detail. The entry builds $A$ from an infinite tower of such reversions: in
the notation of its examples, $A=\mathrm{INV}(x-x\,F_1)$, where $F_1=\mathrm{INV}(x-2x\,F_2)$, where
$F_2=\mathrm{INV}(x-3x\,F_3)$, and so on without end. It is not obvious what such an infinite tower means;
Section~\ref{sec:results} gives it a precise meaning.

The first terms, starting at $n=1$, are 1, 1, 4, 29, 312, 4528, 83736, 1899663, 51453040, \dots{} They grow very quickly.

\subsection{The question}
Hanna's conjecture says that $a(n)$ is odd exactly when $n$ is one of $1,2,4,8,16,\dots$. The table shows the parity of
the first terms ($1$ means odd, $0$ means even). The powers of $2$ are printed in bold.
\begin{center}\small\setlength{\tabcolsep}{3.2pt}
\begin{tabular}{>{\raggedright\arraybackslash}p{6.4em}>{\centering\arraybackslash}p{1.15em}>{\centering\arraybackslash}p{1.15em}>{\centering\arraybackslash}p{1.15em}>{\centering\arraybackslash}p{1.15em}>{\centering\arraybackslash}p{1.15em}>{\centering\arraybackslash}p{1.15em}>{\centering\arraybackslash}p{1.15em}>{\centering\arraybackslash}p{1.15em}>{\centering\arraybackslash}p{1.15em}>{\centering\arraybackslash}p{1.15em}>{\centering\arraybackslash}p{1.15em}>{\centering\arraybackslash}p{1.15em}>{\centering\arraybackslash}p{1.15em}>{\centering\arraybackslash}p{1.15em}>{\centering\arraybackslash}p{1.15em}>{\centering\arraybackslash}p{1.15em}}
\toprule
$n$ & \hl{1} & \hl{2} & 3 & \hl{4} & 5 & 6 & 7 & \hl{8} & 9 & 10 & 11 & 12 & 13 & 14 & 15 & \hl{16} \\ \midrule
$a(n)\bmod 2$ & \hl{1} & \hl{1} & 0 & \hl{1} & 0 & 0 & 0 & \hl{1} & 0 & 0 & 0 & 0 & 0 & 0 & 0 & \hl{1} \\
\bottomrule
\end{tabular}
\\[4pt]
\begin{tabular}{>{\raggedright\arraybackslash}p{6.4em}>{\centering\arraybackslash}p{1.15em}>{\centering\arraybackslash}p{1.15em}>{\centering\arraybackslash}p{1.15em}>{\centering\arraybackslash}p{1.15em}>{\centering\arraybackslash}p{1.15em}>{\centering\arraybackslash}p{1.15em}>{\centering\arraybackslash}p{1.15em}>{\centering\arraybackslash}p{1.15em}>{\centering\arraybackslash}p{1.15em}>{\centering\arraybackslash}p{1.15em}>{\centering\arraybackslash}p{1.15em}>{\centering\arraybackslash}p{1.15em}>{\centering\arraybackslash}p{1.15em}>{\centering\arraybackslash}p{1.15em}>{\centering\arraybackslash}p{1.15em}>{\centering\arraybackslash}p{1.15em}}
\toprule
$n$ & 17 & 18 & 19 & 20 & 21 & 22 & 23 & 24 & 25 & 26 & 27 & 28 & 29 & 30 & 31 & \hl{32} \\ \midrule
$a(n)\bmod 2$ & 0 & 0 & 0 & 0 & 0 & 0 & 0 & 0 & 0 & 0 & 0 & 0 & 0 & 0 & 0 & \hl{1} \\
\bottomrule
\end{tabular}
\\[4pt]
\begin{tabular}{>{\raggedright\arraybackslash}p{6.4em}>{\centering\arraybackslash}p{1.15em}>{\centering\arraybackslash}p{1.15em}>{\centering\arraybackslash}p{1.15em}>{\centering\arraybackslash}p{1.15em}>{\centering\arraybackslash}p{1.15em}>{\centering\arraybackslash}p{1.15em}>{\centering\arraybackslash}p{1.15em}>{\centering\arraybackslash}p{1.15em}>{\centering\arraybackslash}p{1.15em}>{\centering\arraybackslash}p{1.15em}>{\centering\arraybackslash}p{1.15em}>{\centering\arraybackslash}p{1.15em}>{\centering\arraybackslash}p{1.15em}>{\centering\arraybackslash}p{1.15em}>{\centering\arraybackslash}p{1.15em}>{\centering\arraybackslash}p{1.15em}}
\toprule
$n$ & 33 & 34 & 35 & 36 & 37 & 38 & 39 & 40 & 41 & 42 & 43 & 44 & 45 & 46 & 47 & 48 \\ \midrule
$a(n)\bmod 2$ & 0 & 0 & 0 & 0 & 0 & 0 & 0 & 0 & 0 & 0 & 0 & 0 & 0 & 0 & 0 & 0 \\
\bottomrule
\end{tabular}
\end{center}


\subsection{What this paper proves}
\textbf{Main results.} \emph{(1) The infinite tower of reversions defines exactly one series, and cutting it off after
$d$ levels gives the correct coefficients up to degree $d+1$; this holds for arbitrary weights in place of
$1,2,3,\dots$ (Theorem~\ref{thm:tails}). (2) The coefficient $a(n)$ is odd if and only if $n$ is a power of $2$
(Theorem~\ref{thm:parity}).}

The reason for (2) is that, modulo $2$, the weight $2$ at the second level of the tower vanishes. This cuts the tower off
after its first level, and what remains is the equation $A=x+A^2$, which defines the Catalan numbers. Their behaviour
modulo $2$ is classical.

\subsection{What was known}
The sequence was added to the OEIS in 2011, and the conjecture in April 2026. The entry gives no proof. Its computer
program evaluates a tower of finite height, which Theorem~\ref{thm:tails} shows is correct. The parity of the Catalan
numbers is classical; see for example~\cite{deutschsagan}. We are not aware of an earlier proof of the conjecture.

\subsection{Why it is interesting}
Infinitely nested expressions are a common way of describing power series in the OEIS, and they raise a natural
question: what does the expression mean? Our first theorem answers this for every tower of reversions of this shape: each
additional level fixes one more coefficient, so the coefficients of the truncated towers eventually stop changing. The
parity proof then shows a pleasant phenomenon: working modulo a number $p$ can turn an infinite object into a finite one,
because the weight $p$ in the tower becomes $0$. Section~\ref{sec:discussion} explains what this gives for other primes
$p$.

\subsection{How to read this paper}
Section~\ref{sec:background} explains the notions we need: power series, substitution and reversion, remainders, and the
Catalan numbers modulo $2$. Section~\ref{sec:results} states the results precisely, and Section~\ref{sec:proofs} proves
them. Section~\ref{sec:checks} describes how the results were checked, by computer and by a formal proof in Lean.
Section~\ref{sec:discussion} discusses the scope of the result and some open questions.
\section{Background}\label{sec:background}
This section collects the definitions and facts used later. Readers who know formal power series and congruences can skip it and return when needed.
\subsection{Power series as bookkeeping}
A \emph{formal power series} is a way of storing an infinite sequence of numbers $u_0,u_1,u_2,\dots$ in one object,
\[
U(x)=u_0+u_1x+u_2x^2+u_3x^3+\cdots .
\]
The letter $x$ is only a placeholder that keeps the terms apart. We never substitute a number for $x$, so the question
of whether the series ``converges'' never arises; a power series is simply its list of coefficients. We write
$\coeff{n}{U}=u_n$ for the coefficient of $x^n$ in $U$. Two series are equal when all their coefficients are equal.

Series are added coefficient by coefficient and multiplied like polynomials:
\[
\coeff{n}{(UV)}=u_0v_n+u_1v_{n-1}+\cdots+u_nv_0 .
\]
Each coefficient of a product is a finite sum, so everything is well defined. For example,
$(1-x)(1+x+x^2+x^3+\cdots)=1$, because all the other terms cancel. So $1+x+x^2+\cdots$ is the
\emph{multiplicative inverse} of $1-x$, and we write it as $\frac{1}{1-x}$. In the same way, every series $D$ whose
constant term is $1$ has exactly one multiplicative inverse $\frac1D$. Its coefficients can be computed one at a time
from $D\cdot\frac1D=1$, using only additions, subtractions and multiplications. Every fraction in this paper has a
denominator of this kind, so it stands for an honest power series.

Two more pieces of notation. We write $U=O(x^d)$ when all coefficients of $U$ below degree $d$ are zero. So
$U-V=O(x^d)$ means that $U$ and $V$ agree up to degree $d-1$. A series of the form $x+u_2x^2+u_3x^3+\cdots$, with
constant term $0$ and coefficient of $x$ equal to $1$, is called \emph{normalized}.

\subsection{Substituting one series into another}
If $V(x)$ has constant term $0$, we can substitute $V$ for $x$ in $U$:
\[
U(V(x))=u_0+u_1V(x)+u_2V(x)^2+u_3V(x)^3+\cdots .
\]
This is called \emph{composition}, also written $U\circ V$. It makes sense because $V(x)^k$ has no terms below $x^k$,
so only the finitely many terms with $k\le n$ affect the coefficient of $x^n$. For example, if $U=x+x^2$ and $V=x^2$, then
$U(V(x))=x^2+x^4$. Composition is associative: $U\circ(V\circ W)=(U\circ V)\circ W$.

Substituting a series into itself repeatedly gives its \emph{iterates}: $U^{\circ2}(x)=U(U(x))$,
$U^{\circ3}(x)=U(U(U(x)))$, and so on, with $U^{\circ0}(x)=x$ and $U^{\circ1}=U$. They must not be confused with
\emph{powers} $U^2=U\cdot U$, $U^3=U\cdot U\cdot U$. (The OEIS sometimes writes $A^k(x)$ for the $k$-th iterate. In
this paper the small circle always means an iterate, and a plain exponent always means a power.)

\subsection{Undoing a substitution}
A normalized series can always be ``undone'' by another normalized series.

\begin{lemma}[Compositional inverse]\label{lem:inverse}
For every normalized series $U$ there is exactly one normalized series $V$ with $U(V(x))=x$, and this $V$ also
satisfies $V(U(x))=x$. If $U$ has integer coefficients, so does $V$.
\end{lemma}

We call $V$ the \emph{compositional inverse} of $U$. The OEIS calls it the \emph{series reversion} of $U$. For example,
the compositional inverse of $x-x^2$ is $C(x)=x+x^2+2x^3+5x^4+14x^5+\cdots$, whose coefficients are the Catalan numbers.
Indeed $C-C^2=x$, which says exactly that substituting $C$ into $x-x^2$ gives $x$.

\begin{proof}
Write $V=x+v_2x^2+v_3x^3+\cdots$ with unknown coefficients. In
$U(V)=V+u_2V^2+u_3V^3+\cdots$, the coefficient of $x^n$ (for $n\ge2$) equals $v_n$ plus terms that involve only
$v_2,\dots,v_{n-1}$. Setting it equal to $0$ therefore determines $v_n$ from the earlier coefficients, using only
additions, subtractions and multiplications. So there is exactly one $V$ with $U\circ V=x$, and it has integer
coefficients if $U$ does. Applying the same construction to $V$ gives a normalized $W$ with $V\circ W=x$. Then
$U=U\circ(V\circ W)=(U\circ V)\circ W=W$, so $V\circ U=V\circ W=x$ as well.
\end{proof}

\subsection{Remainders and congruences}
For an integer $m\ge2$ we write $a\equiv b\pmod m$ (``$a$ is congruent to $b$ modulo $m$'') when $a-b$ is divisible by
$m$, that is, when $a$ and $b$ leave the same remainder after division by $m$. For example, $17\equiv5\pmod 4$, since both leave remainder $1$, and $-3\equiv1\pmod4$. The
remainder of $a$ on division by $m$ is called its \emph{residue} modulo $m$; it is one of $0,1,\dots,m-1$. Saying that
$a$ is odd is the same as saying $a\equiv1\pmod2$.

Residues can be added and multiplied consistently, as on a clock. The residues $0,1,\dots,m-1$ with this ``clock
arithmetic'' form a number system written $\Z/m\Z$. When $m=p$ is a prime we also write $\F_p$. In $\F_2=\{0,1\}$ we
have $1+1=0$, so adding and subtracting are the same thing.

For power series with integer coefficients, $U\equiv V\pmod m$ means that \emph{every} coefficient of $U-V$ is
divisible by $m$. Equivalently, the two series become equal when every coefficient is replaced by its residue; we call
this \emph{reducing} the series modulo $m$. The key fact is simple:

\emph{Reduction modulo $m$ is compatible with all the operations above.} To reduce a sum, a product, a substitution, or
an inverse, we may reduce the ingredients first and then compute with clock arithmetic. This holds because every
coefficient of the result is obtained from finitely many coefficients of the ingredients by adding, subtracting and
multiplying, and these operations are compatible with remainders.

\para{A word on ``rings''.} Several lemmas below hold not only for integers, but in every number system in
which one can add, subtract and multiply with the usual rules (such as $\Z$, $\Q$, and $\Z/m\Z$). Mathematicians call
such a number system a \emph{commutative ring}. We use the word only in this sense. No division is ever needed, which
matters because in $\Z/m\Z$ one cannot always divide (for instance, $2$ has no inverse in $\Z/4\Z$).

\subsection{The method in one paragraph}
Almost every proof in this paper follows the same three steps.
\begin{enumerate}
\item \emph{Uniqueness.} Show that the defining equation determines the coefficients one at a time, using only
addition, subtraction and multiplication. Then the equation has exactly one solution of the required shape, in every
commutative ring, including the clock arithmetic $\Z/2\Z$.
\item \emph{Reduction.} Reducing the integer solution modulo $2$ gives a solution of the reduced equation,
because reduction is compatible with all operations.
\item \emph{A good guess.} Find a simple series that satisfies the reduced equation and check that it does. By uniqueness,
this simple series \emph{is} the reduction of the sequence, and its coefficients can simply be read off.
\end{enumerate}
The creative part is step~3: finding a simple series that works. Step~1 sometimes needs care, when the equation seems to
require a division.
For a first example, take the Catalan series $C=x+C^2$ from above. Its coefficients $1,1,2,5,14,42,\dots$ are
determined one at a time by this equation. Modulo $2$, the series $S=x+x^2+x^4+x^8+\cdots$ satisfies the same
equation: in $\F_2$, $S^2=x^2+x^4+x^8+\cdots$ (see Lemma~\ref{lem:frobenius}), so $x+S^2=S$. Hence $C\equiv S\pmod2$,
and the Catalan coefficient $\coeff{n}{C}$ is odd exactly when $n$ is a power of $2$.

\subsection{Squaring in clock arithmetic modulo \texorpdfstring{$2$}{2}}
\begin{lemma}[Frobenius identity]\label{lem:frobenius}
Let $p$ be a prime. For every power series $U$ with coefficients in $\F_p$, we have $U(x)^p=U(x^p)$.
\end{lemma}

\begin{proof}
In $\F_p$, the binomial theorem simplifies to $(a+b)^p=a^p+b^p$, because every binomial coefficient
$\binom pi$ with $0<i<p$ is divisible by $p$. Also $c^p=c$ for every $c\in\F_p$ (Fermat's little theorem). So
$\bigl(\sum_nu_nx^n\bigr)^p=\sum_nu_n^px^{pn}=\sum_nu_nx^{pn}=U(x^p)$. Only finitely many terms matter for any given
coefficient, so this is a polynomial identity.
\end{proof}

For $p=2$ this says that squaring a series modulo $2$ just doubles every exponent: for instance,
$(x+x^3)^2\equiv x^2+x^6\pmod2$. In particular, the square of any series has, modulo $2$, no terms of odd degree.

\subsection{The Catalan numbers modulo 2}
\begin{lemma}\label{lem:catalan}
In every commutative ring the equation $U=x+U^2$ has exactly one solution with $U(0)=0$. Over the integers it is the
Catalan series $C(x)=x+x^2+2x^3+5x^4+14x^5+\cdots$, and
\[
C(x)\equiv x+x^2+x^4+x^8+\cdots=\sum_{k\ge0}x^{2^k}\pmod 2 .
\]
So $\coeff{n}{C}$ is odd exactly when $n$ is a power of $2$.
\end{lemma}

\begin{proof}
The coefficient of $x^n$ in $U^2$ involves only $u_1,\dots,u_{n-1}$, so the equation determines $u_n$ one at a time:
$u_1=1$ and $u_n=\sum_{i=1}^{n-1}u_iu_{n-i}$. Over $\F_2$, the series $S=\sum_kx^{2^k}$ satisfies
$x+S^2=x+\sum_kx^{2^{k+1}}=S$ by Lemma~\ref{lem:frobenius}, so by uniqueness $C\equiv S\pmod2$.
\end{proof}
The parity of the Catalan numbers is classical; see for example~\cite{deutschsagan}. (Here the Catalan number
$\frac1{n}\binom{2n-2}{n-1}$ appears as the coefficient of $x^n$, so ``$n$ a power of $2$'' corresponds to the familiar
statement that the $k$-th Catalan number is odd exactly when $k+1$ is a power of $2$.)


\section{Results}\label{sec:results}

Following the OEIS, we write $\mathrm{INV}(F)$ for the compositional inverse of a normalized series $F$
(Lemma~\ref{lem:inverse}).

\para{Towers of reversions.} Let $R$ be a commutative ring, and let $w_0,w_1,w_2,\dots$ be elements of $R$, called the
\emph{weights}. A \emph{solution of the tower} is a sequence of normalized series $F_0,F_1,F_2,\dots$ over $R$ such that
\begin{equation}\label{eq:tails}
F_j=\mathrm{INV}\bigl(x-w_j\,x\,F_{j+1}\bigr)\qquad\text{for every }j\ge0 .
\end{equation}
Then $F_0=\mathrm{INV}(x-w_0xF_1)=\mathrm{INV}\bigl(x-w_0x\,\mathrm{INV}(x-w_1xF_2)\bigr)=\cdots$. For OEIS A195194 the
weights are $w_j=j+1$, so $F_0$ is the series in the definition of the entry. The \emph{tower of height $d$} is obtained
by replacing the innermost $\mathrm{INV}(\cdots)$ by $x$ after $d$ levels: we define $F^{[d]}_j$ for all $j$ by
\[
F^{[0]}_j=x,\qquad F^{[d+1]}_j=\mathrm{INV}\bigl(x-w_j\,x\,F^{[d]}_{j+1}\bigr).
\]
For example, $F^{[2]}_0=\mathrm{INV}\bigl(x-w_0x\,\mathrm{INV}(x-w_1x^2)\bigr)$.

\begin{theorem}[The tower is well defined]\label{thm:tails}
For every commutative ring $R$ and all weights $w_0,w_1,\dots$ in $R$, the tower \eqref{eq:tails} has exactly one
solution. Moreover, the towers of finite height approach it: $F^{[d]}_j-F_j=O(x^{d+2})$ for all $d,j\ge0$. That is,
$F^{[d]}_j$ and $F_j$ have the same coefficients up to degree $d+1$.
\end{theorem}

So the infinite tower in the definition of the entry has exactly one meaning, and the sequence A195194 is
$a(n)=\coeff{n}{F_0}$ for the weights $w_j=j+1$ over the integers.

\begin{theorem}[Hanna's conjecture]\label{thm:parity}
For the weights $w_j=j+1$ over the integers, the coefficient $a(n)=\coeff{n}{F_0}$ is odd if and only if $n$ is a power
of $2$. More precisely,
\[
F_0(x)\equiv x+x^2+x^4+x^8+\cdots=\sum_{k\ge0}x^{2^k}\pmod 2 .
\]
\end{theorem}

\section{Proofs}\label{sec:proofs}

\para{The idea in brief.} Each reversion in the tower can be rewritten as an ordinary equation:
$F_j=\mathrm{INV}(x-w_jxF_{j+1})$ says that $F_j=x+w_j\,F_j\cdot F_{j+1}(F_j)$. The right-hand side, in degree $n$,
only involves coefficients of degree less than $n$. So all the series $F_j$ can be built at the same time, one degree at
a time; this gives existence and uniqueness. Modulo $2$, the weight $w_1=2$ is $0$, which forces $F_1=x$; then the
equation for $F_0$ becomes the Catalan equation $F_0=x+F_0^2$.

\begin{lemma}\label{lem:agree}
Let $U$ and $U'$ be normalized series that agree below degree $r$, that is, $U-U'=O(x^r)$. Then their compositional
inverses $V=\mathrm{INV}(U)$ and $V'=\mathrm{INV}(U')$ also agree below degree $r$.
\end{lemma}

\begin{proof}
Since $U(V)=x=U'(V')$, we have $U(V)-U(V')=U'(V')-U(V')=(U'-U)(V')$. The series $U'-U$ has no terms below degree $r$,
and $V'$ starts with $x$, so $(U'-U)(V')=O(x^r)$. Now suppose that $V-V'$ were not $O(x^r)$, and let $s<r$ be the lowest
degree in which $V$ and $V'$ differ. Then
\[
U(V)-U(V')=(V-V')+\sum_{k\ge2}u_k\,\bigl(V^k-V'^k\bigr).
\]
Each $V^k-V'^k=(V-V')(V^{k-1}+V^{k-2}V'+\cdots+V'^{k-1})$ has no terms below degree $s+k-1\ge s+1$. So the right-hand
side has a nonzero coefficient in degree $s$, coming from $V-V'$. This contradicts $U(V)-U(V')=O(x^r)$.
\end{proof}

\begin{proof}[Proof of Theorem~\ref{thm:tails}]
\emph{Rewriting the tower.} Let $F_j$ and $F_{j+1}$ be normalized. The series $x-w_jxF_{j+1}$ is normalized too. By
Lemma~\ref{lem:inverse}, $F_j=\mathrm{INV}(x-w_jxF_{j+1})$ holds if and only if substituting $F_j$ into
$x-w_jxF_{j+1}$ gives $x$, that is, $F_j-w_j\,F_j\cdot F_{j+1}(F_j)=x$, or
\begin{equation}\label{eq:fixed}
F_j=x+w_j\,F_j\cdot F_{j+1}(F_j).
\end{equation}

\emph{Existence and uniqueness.} Both factors of the product $F_j\cdot F_{j+1}(F_j)$ have no constant term. So the
coefficient of $x^n$ in this product involves only coefficients of $F_j$ and $F_{j+1}$ of degree less than $n$. Hence
the whole sequence $F_0,F_1,F_2,\dots$ can be built at the same time, one degree at a time: every $F_j$ has constant term
$0$ and coefficient $1$ in degree $1$; and once all the $F_j$ are known below degree $n$, equation \eqref{eq:fixed}
determines the coefficient of $x^n$ in every $F_j$. This shows that there is exactly one solution. Only additions and
multiplications are used, so this works in every commutative ring.

\emph{Towers of finite height.} We show $F^{[d]}_j-F_j=O(x^{d+2})$ by induction on $d$, for all $j$ at once. For $d=0$,
both $F^{[0]}_j=x$ and $F_j$ are normalized, so they agree below degree $2$. Suppose $F^{[d]}_{j+1}-F_{j+1}=O(x^{d+2})$.
Then the series $x-w_jxF^{[d]}_{j+1}$ and $x-w_jxF_{j+1}$ differ by $w_jx\,(F_{j+1}-F^{[d]}_{j+1})=O(x^{d+3})$. By
Lemma~\ref{lem:agree}, their inverses $F^{[d+1]}_j$ and $F_j$ also differ by $O(x^{d+3})$.
\end{proof}

\begin{proof}[Proof of Theorem~\ref{thm:parity}]
Let $F_0,F_1,\dots$ be the solution over the integers with $w_j=j+1$. Reducing all coefficients modulo $2$ gives
normalized series over $\F_2$ that satisfy \eqref{eq:fixed} with the weights reduced modulo $2$, because reduction is
compatible with all the operations involved. By the uniqueness in Theorem~\ref{thm:tails} for the ring $\F_2$, these
reductions are \emph{the} solution of the tower over $\F_2$. We now work in $\F_2$.

The weight $w_1=2$ is $0$ in $\F_2$, so \eqref{eq:fixed} for $j=1$ says simply $F_1=x$. Then \eqref{eq:fixed} for $j=0$,
with $w_0=1$, says
\[
F_0=x+F_0\cdot F_1(F_0)=x+F_0\cdot F_0=x+F_0^2 .
\]
This is the Catalan equation, and by Lemma~\ref{lem:catalan} its solution over $\F_2$ is $\sum_{k\ge0}x^{2^k}$. So
$a(n)$ is odd exactly when $n$ is a power of $2$.
\end{proof}

Notice that, modulo $2$, all the deeper levels $F_2,F_3,\dots$ of the tower play no role at all.

\section{Checking the results}\label{sec:checks}

A mathematical proof can contain mistakes, so the results were checked in two independent ways.

\para{Numerical check.}
We evaluated the tower with exact integer arithmetic, once with height $143$ and once with height $153$. The two
results agree below degree $140$, as Theorem~\ref{thm:tails} predicts. This reproduces all 20 terms printed in the OEIS
entry, and the parity rule of Theorem~\ref{thm:parity} holds for every $n$ with $1\le n<140$.
The computation was done with the PARI/GP computer algebra system, using code that was written separately from the
formal proofs (file \texttt{verification/independent\_checks.gp} in the accompanying repository~\cite{repo}).
A computation can only test finitely many cases, so it is a check on the definitions and on the indexing, not a proof.

\para{Formal proof in Lean.}
Lean~\cite{lean} is a \emph{proof assistant}: a program in which mathematical statements and their proofs are written
in a precise formal language, and every step is checked mechanically down to the basic rules of logic. All results of
this paper are proved in Lean~4 (version 4.19.0) on top of Mathlib~\cite{mathlib}, Lean's library of formalized
mathematics. The statements are collected in the file \texttt{lean/Entries/A195194.lean} of the repository,
in the namespace \texttt{FinalCollection.A195194}. There the solution of the tower is called \texttt{tail} and the towers of finite height \texttt{approximation}; the proofs are in the module \texttt{Round4/NestedReversion.lean}.
The table lists where each result of this paper is proved.
\begin{center}\small\renewcommand{\arraystretch}{1.25}
\begin{tabular}{>{\raggedright\arraybackslash}p{0.52\linewidth}>{\raggedright\arraybackslash}p{0.4\linewidth}}
\toprule
Statement in this paper & Lean theorem(s)\\
\midrule
Theorem~\ref{thm:tails} for arbitrary weights in any commutative ring & \texttt{arbitrary\_\allowbreak{}weights\_\allowbreak{}unique}, \texttt{arbitrary\_\allowbreak{}finite\_\allowbreak{}depth} \\
For $w_j=j+1$: every $F_j$ is normalized, satisfies both inverse identities, and the solution is unique & \texttt{all\_\allowbreak{}tails\_\allowbreak{}normalized}, \texttt{all\_\allowbreak{}tails\_\allowbreak{}source}, \texttt{nested\_\allowbreak{}definition\_\allowbreak{}unique} \\
$A=F_0$ is normalized and is the inverse of $x-xF_1$ & \texttt{normalized\_\allowbreak{}source}, \texttt{first\_\allowbreak{}tail\_\allowbreak{}reversion} \\
The tower of height $d$ is correct up to degree $d+1$ & \texttt{finite\_\allowbreak{}nest\_\allowbreak{}coefficients} \\
Theorem~\ref{thm:parity}: $F_0\equiv x+F_0^2\pmod2$, and the parity rule & \texttt{binary\_\allowbreak{}equation}, \texttt{exact\_\allowbreak{}parity} \\
\bottomrule
\end{tabular}
\end{center}
The Lean project was rebuilt from source for this paper, and Lean was asked to list the basic assumptions (axioms)
that each of these theorems ultimately depends on. Every one depends only on the three standard axioms of Lean's logic
(\texttt{propext}, \texttt{Classical.choice}, \texttt{Quot.sound}); no step is left unproved. A formal proof certifies
exactly the statement that Lean checks. Whether that statement says what the OEIS entry means is a separate question,
so the Lean definitions were also compared, side by side, with the definition in the OEIS entry
(Section~\ref{sec:sequence}).

\section{Discussion}\label{sec:discussion}

\para{Scope.} Theorem~\ref{thm:parity} determines the parity of every term, so it settles the conjecture of the entry
completely. It says nothing about the size of $a(n)$, which grows very fast, or about remainders modulo other numbers.

\para{Other primes.} The argument of Theorem~\ref{thm:parity} has a general version. Let $p$ be a prime. Modulo $p$, the
weight $w_{p-1}=p$ vanishes, so $F_{p-1}\equiv x$ and the tower becomes finite:
\[
F_0\equiv\mathrm{INV}\Bigl(x-x\,\mathrm{INV}\bigl(x-2x\,\mathrm{INV}(\cdots\mathrm{INV}(x-(p-1)x^2)\cdots)\bigr)\Bigr)
\pmod p .
\]
A power series $F$ over $\F_p$ is called \emph{algebraic} if it satisfies a polynomial equation $Q(x,F(x))=0$, where $Q$
is a nonzero polynomial in two variables. The innermost series $\mathrm{INV}(x-(p-1)x^2)$ is algebraic; and if $U$ is
algebraic, then so is $x-wxU$ for every nonzero weight $w$, and so is $\mathrm{INV}(U)$ (if $Q(x,U(x))=0$, then
substituting $V=\mathrm{INV}(U)$ gives $Q(V(x),x)=0$). So $F_0$ is algebraic modulo $p$. By a theorem of Christol, Kamae,
Mend\`es France and Rauzy~\cite{christol,allouche}, the remainders $a(n)\bmod p$ then form a \emph{$p$-automatic}
sequence: they can be computed by a machine with finitely many states that reads the digits of $n$ in base $p$. For
$p=2$ the machine is very simple: it checks whether $n$ is a power of $2$. We have not worked out the machines for
$p\ge3$, and this seems a natural question. (This paragraph is not part of the Lean formalization.)

\para{The method.} The proof illustrates a principle used throughout this collection of papers: a sequence defined by an
equation between power series often becomes simple after reduction modulo a small number. Here the reduction even turns
an infinite construction into a finite one.

\section*{About this paper}\label{sec:about}
\para{How the work was done.}
This paper is one of 34 papers on conjectures from the OEIS that came out of an autonomous AI research effort directed
by Dev~D.~Goyal~\cite{repo}. The roles were as follows.
\begin{itemize}
\item \emph{Dev~D.~Goyal} started the project, chose which problems to pursue, asked for proofs, formal verification and
revisions, and decided what to publish. Dev~D.~Goyal did not supply proof strategies or intermediate arguments.
\item \emph{ChatGPT} (OpenAI), using the GPT-6 model with Pro reasoning, worked autonomously from those prompts. It
searched the OEIS for suitable conjectures, found the proofs, wrote the Lean~4 formalization, and wrote a first draft
of this paper.
\item \emph{Claude Opus~5.5}, used with Claude Code, checked the work independently and then rewrote this paper. It
rebuilt the Lean project from source and confirmed that the theorems listed in Section~\ref{sec:checks} rely only on
Lean's standard axioms; compared the Lean definitions with the OEIS definition; recomputed the sequence with separately
written PARI/GP code, compared it with the terms published in the OEIS, and tested every claim of the paper on many
terms; and read every step of the proofs. It then rewrote the text to make it self-contained and readable. The
theorems and the proof strategies are ChatGPT's. In rewriting, Claude Opus~5.5 expanded the explanations and, in
places, wrote out an argument in more detail or by a more elementary route. The Discussion section was written during
this rewrite; the further observations it contains were checked by computer, and they are not part of the Lean
formalization.
\end{itemize}
Dev~D.~Goyal has verified the results and is responsible for this paper.

\para{Review status.} This paper has not been peer reviewed, and it has not been submitted to a journal or to the OEIS.

\para{Priority.} An OEIS entry that labels a statement as a conjecture shows only that its contributor did not know
a proof; it does not show that no proof existed. During the research only limited searches for earlier work were
made. No claim of priority is made.

\para{Corrections.} Please report any error in the issue tracker of the repository~\cite{repo}, where the Lean
sources, the numerical checks, and the earlier drafts are also available.

\section*{Acknowledgements}
The sequence and the conjecture are due to Paul~D.~Hanna~\cite{oeis}.

\begin{thebibliography}{99}\small
\raggedright
\bibitem{oeis} P.~D.~Hanna, entry A195194 (G.f.: $A(x)=\mathrm{INV}(x-x\,\mathrm{INV}(x-2x\,\mathrm{INV}(x-3x\,\mathrm{INV}(\cdots))))$), \emph{The On-Line Encyclopedia of Integer Sequences}, 2011; conjecture April 2026. \url{https://oeis.org/A195194}. Text as of October 2026.
\bibitem{deutschsagan} E.~Deutsch and B.~E.~Sagan, Congruences for Catalan and Motzkin numbers and related sequences, \emph{Journal of Number Theory} 117 (2006), 191--215.
\bibitem{christol} G.~Christol, T.~Kamae, M.~Mend\`es France and G.~Rauzy, Suites alg\'ebriques, automates et substitutions, \emph{Bulletin de la Soci\'et\'e Math\'ematique de France} 108 (1980), 401--419.
\bibitem{allouche} J.-P.~Allouche and J.~Shallit, \emph{Automatic Sequences}, Cambridge University Press, 2003.
\bibitem{oeisdb} N.~J.~A.~Sloane and The OEIS Foundation Inc., \emph{The On-Line Encyclopedia of Integer Sequences}. \url{https://oeis.org}.
\bibitem{stanley} R.~P.~Stanley, \emph{Enumerative Combinatorics}, Volume~2, Cambridge University Press, 1999. (Chapter~5 treats formal power series, composition and compositional inverses.)
\bibitem{niven} I.~Niven, Formal power series, \emph{American Mathematical Monthly} 76 (1969), 871--889.
\bibitem{lean} L.~de~Moura and S.~Ullrich, The Lean~4 theorem prover and programming language, in \emph{Automated Deduction -- CADE 28}, Lecture Notes in Computer Science 12699, Springer, 2021, 625--635.
\bibitem{mathlib} The mathlib Community, The Lean mathematical library, in \emph{Proceedings of the 9th ACM SIGPLAN International Conference on Certified Programs and Proofs (CPP 2020)}, 367--381. Version used: commit \texttt{c44e0c8ee63c}.
\bibitem{repo} D.~D.~Goyal, \emph{Proofs and corrections for 34 OEIS entries, with Lean~4 formalizations}, repository, 2026. \url{https://github.com/dvgyl/oeis-a396843-proof}.
\end{thebibliography}
\end{document}
